\documentclass[11pt]{article}

% Required Packages
\usepackage{amsmath}
\usepackage{amssymb}
\usepackage{amsthm}
\usepackage{geometry}
\usepackage[colorlinks=true, linkcolor=blue, bookmarks=true]{hyperref}
\usepackage{mathrsfs}
\usepackage{mathtools}
\usepackage{tikz}

% Page Geometry
\geometry{letterpaper, margin=1in}

% Theorem Environments
\newtheorem{theorem}{Theorem}
\newtheorem{lemma}{Lemma}
\newtheorem{proposition}{Proposition}
\newtheorem{corollary}{Corollary}
\theoremstyle{definition}
\newtheorem{problem}{Problem}
\newtheorem{definition}{Definition}

% Title Information
\title{\textbf{Advanced Mathematics: Existence, Logical Boundaries, and Proof by Contradiction} \\ \Large Comprehensive Practice Problem Set}
\author{Tashfeen Omran}
\date{\today}

\begin{document}

\maketitle

\tableofcontents
\newpage

% ==========================================
% SECTION 1: CONCEPT CHECKLIST
% ==========================================
\section{Concept Checklist \& Question Mapping}
\label{sec:checklist_start}

This problem set is rigorously derived from the theoretical frameworks of logical boundary testing, universal falsification, and impossibility proofs. Every problem maps directly to the core concepts extracted from the source material.

\begin{itemize}
    \item \textbf{C1: Universal Falsification (Finite Domains):} Exhaustive testing to find edge cases where the hypothesis holds but the conclusion fails. \textit{(Mapped to Problems 1, 2)}
    \item \textbf{C2: Universal Falsification (Infinite Domains):} Analyzing algebraic structures, inequalities, and domain restrictions (e.g., logarithms, roots) to construct counterexamples. \textit{(Mapped to Problems 3, 4, 5)}
    \item \textbf{C3: Impossibility Proofs (Parity Forcing):} Mapping complex polynomial expressions to the binary space of $\{0, 1\} \pmod{2}$ to expose contradictions. \textit{(Mapped to Problems 6, 7, 8)}
    \item \textbf{C4: Impossibility Proofs (Divisibility \& Primes):} Forcing violations of the fundamental theorem of arithmetic or basic divisibility axioms. \textit{(Mapped to Problems 9, 10, 11)}
    \item \textbf{C5: Rationality Bounds \& Density:} Exploiting the density of $\mathbb{Q}$ and $\mathbb{R} \setminus \mathbb{Q}$ using the ``Divide by 2'' trick to shatter extremum assumptions. \textit{(Mapped to Problems 12, 13, 14)}
    \item \textbf{C6: De Morgan's Laws on Implications:} Negating disjunctive conclusions ($\sim(P \implies Q \lor R) \equiv P \land \sim Q \land \sim R$) to form concrete simultaneous equations. \textit{(Mapped to Problems 15, 16)}
    \item \textbf{C7: Fermat's Method of Infinite Descent:} Constructing an infinite strictly decreasing sequence of positive integers to violate the Well-Ordering Principle of $\mathbb{N}$. \textit{(Mapped to Problems 17, 18, 19)}
    \item \textbf{C8: Set-Theoretic \& Probability Foundations:} Utilizing mutually exclusive/exhaustive sets and probability measure axioms ($\mathbb{P}(A)=1$) to prove logical necessities. \textit{(Mapped to Problems 20, 21, 22)}
    \item \textbf{C9: Abstract Real-World Modeling (Quant Finance):} Constructing self-financing portfolios to contradict the existence of an Equivalent Martingale Measure ($\mathbb{Q}$) in Arbitrage Pricing Theory. \textit{(Mapped to Problems 23, 24, 25)}
    \item \textbf{C10: Abstract Real-World Modeling (Machine Learning):} Proving PAC learning lower bounds by assuming learnability for infinite VC-dimension classes and deriving contradictions. \textit{(Mapped to Problems 26)}
    \item \textbf{C11: Advanced Diagnostic Testing (Find the Flaw):} Identifying subtle logical fallacies, such as unproven conditionals, division by zero, or invalid theorem applications (e.g., IVT on discontinuous functions). \textit{(Mapped to Problems 27, 28, 29, 30)}
\end{itemize}

\newpage

% ==========================================
% SECTION 2: PROBLEM SETS
% ==========================================
\section{Problem Sets}
\label{sec:problems}

\subsection{Part A: Universal Falsification (Counterexamples)}

\begin{problem}
Disprove the statement: If $n \in \{1, 2, 3, 4, 5\}$, then $2^n + n^2$ is a prime number.
\end{problem}

\begin{problem}
Disprove the statement: If $n \in \{1, 2, 3, 4, 5\}$, then $4 \mid (3n^2 + 1)$.
\end{problem}

\begin{problem}
Disprove the statement: For every two positive real numbers $x$ and $y$, $\sqrt{x^2 + y^2} = x + y$.
\end{problem}

\begin{problem}
Disprove the statement: For all real numbers $x$, if $x^2 > x$, then $x > 1$.
\end{problem}

\begin{problem}
Disprove the statement: For all real numbers $a, b \in \mathbb{R}$, if $\log(a^2) = \log(b^2)$, then $a = b$.
\end{problem}

\subsection{Part B: Impossibility Proofs (Parity \& Divisibility)}

\begin{problem}
Use proof by contradiction to prove that $500$ cannot be written as the sum of an odd integer and two even integers.
\end{problem}

\begin{problem}
Prove by contradiction that if $a$ and $b$ are odd integers, then $a^2 + b^2$ is not divisible by $4$.
\end{problem}

\begin{problem}
Prove by contradiction that there are no integer solutions to the equation $x^2 - 4y = 2$.
\end{problem}

\begin{problem}
Prove by contradiction that if $a, b \in \mathbb{Z}$ and $a \ge 2$, then $a \nmid b$ or $a \nmid (b + a - 1)$.
\end{problem}

\begin{problem}
Prove by contradiction that for any integer $n$, if $n^2$ is even, then $n$ is even.
\end{problem}

\begin{problem}
Prove by contradiction that the sum of three consecutive integers cannot be equal to $100$.
\end{problem}

\subsection{Part C: Rationality Bounds \& Extremum in Dense Sets}

\begin{problem}
Prove by contradiction that there is no largest negative rational number.
\end{problem}

\begin{problem}
Prove by contradiction that there is no smallest positive irrational number.
\end{problem}

\begin{problem}
Prove by contradiction that the sum of a rational number and an irrational number is always irrational.
\end{problem}

\subsection{Part D: De Morgan's Laws \& Set-Theoretic Logic}

\begin{problem}
Prove by contradiction: For any real numbers $x$ and $y$, if $x + y > 10$, then $x > 5$ or $y > 5$.
\end{problem}

\begin{problem}
Let $A$ and $B$ be sets. Prove by contradiction that if $A \cap B = \emptyset$, then $A \subseteq B^c$.
\end{problem}

\subsection{Part E: Fermat's Method of Infinite Descent}

\begin{problem}
Use the method of infinite descent to prove that $\sqrt{3}$ is irrational.
\end{problem}

\begin{problem}
Use the method of infinite descent to prove that the equation $x^3 + 3y^3 = 9z^3$ has no non-zero integer solutions.
\end{problem}

\begin{problem}
Use the method of infinite descent to prove that there are no positive integer solutions to $x^2 + y^2 = 3z^2$.
\end{problem}

\subsection{Part F: Probability Foundations \& Abstract Modeling}

\begin{problem}
Let $(\Omega, \mathcal{F}, \mathbb{P})$ be a probability space. Prove by contradiction that if $\mathbb{P}(A \cup B) = 1$ and $\mathbb{P}(A) = 0$, then $\mathbb{P}(B) = 1$.
\end{problem}

\begin{problem}
Prove by contradiction: For any two events $A$ and $B$ in a probability space, if $\mathbb{P}(A) + \mathbb{P}(B) > 1$, then $A \cap B \neq \emptyset$.
\end{problem}

\begin{problem}
Prove by contradiction: If an event $A$ is independent of itself, then $\mathbb{P}(A) = 0$ or $\mathbb{P}(A) = 1$.
\end{problem}

\begin{problem}
\textbf{Arbitrage Pricing Theory:} Consider a 1-period financial market with a risk-free bond yielding $r = 0.05$ (so $B_0 = 1, B_1 = 1.05$). A risky asset has initial price $S_0 = 100$. At time $t=1$, the asset price will be either $S_1 = 110$ or $S_1 = 120$. Prove by contradiction that this market contains an arbitrage opportunity.
\end{problem}

\begin{problem}
\textbf{Law of One Price:} In a financial market, let $V_t(A)$ and $V_t(B)$ represent the values of two portfolios at time $t$. Assume that at maturity $T$, $V_T(A) = V_T(B)$ almost surely. Prove by contradiction that if $V_0(A) < V_0(B)$, the market contains an arbitrage opportunity.
\end{problem}

\begin{problem}
\textbf{Equivalent Martingale Measure:} Prove by contradiction that if a market admits an arbitrage opportunity, there cannot exist a risk-neutral probability measure $\mathbb{Q}$ equivalent to the real-world measure $\mathbb{P}$ under which discounted asset prices are martingales.
\end{problem}

\begin{problem}
\textbf{PAC Learning Lower Bounds:} In Statistical Learning Theory, the Fundamental Theorem states that the sample complexity $m$ required to PAC-learn a hypothesis class $\mathcal{H}$ is bounded below by $m \ge C \frac{\text{VC}(\mathcal{H}) + \ln(1/\delta)}{\epsilon}$ for some constant $C > 0$. Prove by contradiction that a hypothesis class with infinite VC-dimension cannot be PAC learned.
\end{problem}

\subsection{Part G: Advanced Diagnostic Testing (Find the Flaw)}

\begin{problem}
\textbf{Find the Flaw:} Identify the logical error in the following proof that $1 = 2$. \\
\textit{Proof.} Assume, to the contrary, that $1 \neq 2$. Let $a = b = 1$. Multiply both sides by $a$: $a^2 = ab$. Subtract $b^2$ from both sides: $a^2 - b^2 = ab - b^2$. Factor both sides: $(a - b)(a + b) = b(a - b)$. Divide both sides by $(a - b)$: $a + b = b$. Substitute $a = 1, b = 1$: $1 + 1 = 1 \implies 2 = 1$. This contradicts our assumption that $1 \neq 2$. Therefore, $1 = 2$. $\blacksquare$
\end{problem}

\begin{problem}
\textbf{Find the Flaw:} Identify the logical error in the following proof that $\sqrt{4}$ is irrational. \\
\textit{Proof.} Assume, to the contrary, that $\sqrt{4}$ is rational. Then $\sqrt{4} = \frac{p}{q}$ where $p, q \in \mathbb{Z}$, $q \neq 0$, and the fraction is in lowest terms. Squaring both sides gives $4 = \frac{p^2}{q^2}$, so $p^2 = 4q^2 = 2(2q^2)$. Since $p^2$ is a multiple of $2$, $p^2$ is even, which implies $p$ is even. Let $p = 2k$. Substitute this back: $(2k)^2 = 4q^2 \implies 4k^2 = 4q^2 \implies k^2 = q^2$. Since $k^2 = q^2$, if $k$ is even, $q$ is even. If $q$ is even, then both $p$ and $q$ are divisible by $2$. This contradicts the assumption that $\frac{p}{q}$ is in lowest terms. Therefore, $\sqrt{4}$ is irrational. $\blacksquare$
\end{problem}

\begin{problem}
\textbf{Find the Flaw:} Identify the logical error in the following proof that the function $f(x) = \frac{1}{x}$ has a root in the interval $[-1, 1]$. \\
\textit{Proof.} Assume, to the contrary, that $f(x)$ has no root in $[-1, 1]$. We evaluate the function at the endpoints of the interval: $f(-1) = \frac{1}{-1} = -1$ and $f(1) = \frac{1}{1} = 1$. Since $f(-1) < 0$ and $f(1) > 0$, by the Intermediate Value Theorem, there must exist some $c \in (-1, 1)$ such that $f(c) = 0$. This means $f(x)$ does have a root, which contradicts our assumption. Therefore, $f(x)$ has a root in $[-1, 1]$. $\blacksquare$
\end{problem}

\begin{problem}
\textbf{Find the Flaw:} Identify the logical error in the following proof that all positive integers are equal. \\
\textit{Proof.} We prove by contradiction. Assume there exist positive integers that are not equal. Let $S$ be the set of all pairs of positive integers $(a, b)$ such that $a \neq b$. By the Well-Ordering Principle, there exists a pair in $S$ with the minimum maximum value. Let this pair be $(x, y)$, and let $M = \max(x, y)$. Consider the pair $(x - 1, y - 1)$. Since $x \neq y$, it follows that $x - 1 \neq y - 1$. Thus $(x - 1, y - 1) \in S$. The maximum of this new pair is $\max(x - 1, y - 1) = M - 1$. Since $M - 1 < M$, we have found a pair in $S$ with a smaller maximum. This contradicts the assumption that $(x, y)$ had the minimum maximum. Therefore, no such pair exists, and all positive integers are equal. $\blacksquare$
\end{problem}

\newpage

% ==========================================
% SECTION 3: DETAILED SOLUTIONS
% ==========================================
\section{Detailed Solutions}
\label{sec:solutions}

\subsection*{Solutions to Part A: Universal Falsification}

\textbf{Solution 1:}
\begin{proof}
To disprove the universal statement $\forall n \in \{1, 2, 3, 4, 5\}, P(n)$, we must find a counterexample where the hypothesis holds but the conclusion fails. Let $P(n) = 2^n + n^2$. We test the elements:
\begin{itemize}
    \item For $n = 1$: $P(1) = 2^1 + 1^2 = 2 + 1 = 3$. (Prime)
    \item For $n = 2$: $P(2) = 2^2 + 2^2 = 4 + 4 = 8$.
\end{itemize}
Notice that $8 = 2 \times 4$. Since $8$ has positive divisors other than $1$ and itself, it is a composite number. Therefore, $n = 2$ is a counterexample that disproves the statement.
\end{proof}

\textbf{Solution 2:}
\begin{proof}
We seek an $n \in \{1, 2, 3, 4, 5\}$ such that $3n^2 + 1$ is not divisible by $4$. Let us evaluate the expression for $n = 2$:
\[ 3(2)^2 + 1 = 3(4) + 1 = 12 + 1 = 13 \]
We must check if $4 \mid 13$. Division of $13$ by $4$ yields $13 = 4(3) + 1$. Since the remainder is $1 \neq 0$, $4$ does not divide $13$. Thus, $n = 2$ is a counterexample.
\end{proof}

\textbf{Solution 3:}
\begin{proof}
We seek positive real numbers $x, y \in \mathbb{R}^+$ such that the equality fails. Let $x = 3$ and $y = 4$. Evaluate the left-hand side (LHS):
\[ \sqrt{x^2 + y^2} = \sqrt{3^2 + 4^2} = \sqrt{9 + 16} = \sqrt{25} = 5 \]
Evaluate the right-hand side (RHS):
\[ x + y = 3 + 4 = 7 \]
Since $5 \neq 7$, the equality does not hold. Thus, $x = 3, y = 4$ is a counterexample.
\end{proof}

\textbf{Solution 4:}
\begin{proof}
We must find a real number $x$ such that the hypothesis ($x^2 > x$) is true, but the conclusion ($x > 1$) is false. Let $x = -1$. Evaluate the hypothesis:
\[ (-1)^2 = 1 \]
Since $1 > -1$, the hypothesis $x^2 > x$ holds. Now evaluate the conclusion:
\[ -1 > 1 \quad \text{(False)} \]
Since the hypothesis is true and the conclusion is false, $x = -1$ is a valid counterexample.
\end{proof}

\textbf{Solution 5:}
\begin{proof}
We must find real numbers $a, b$ such that $\log(a^2) = \log(b^2)$ is true, but $a = b$ is false. Let $a = 2$ and $b = -2$. Evaluate the hypothesis:
\[ \log(a^2) = \log(2^2) = \log(4) \]
\[ \log(b^2) = \log((-2)^2) = \log(4) \]
The hypothesis holds since $\log(4) = \log(4)$. However, $2 \neq -2$, so the conclusion $a = b$ fails. Thus, $a = 2, b = -2$ is a counterexample.
\end{proof}

\subsection*{Solutions to Part B: Impossibility Proofs}

\textbf{Solution 6:}
\begin{proof}
We proceed by contradiction. Assume, to the contrary, that $500$ can be written as the sum of an odd integer and two even integers. Let the odd integer be $o = 2k + 1$ for some $k \in \mathbb{Z}$. Let the two even integers be $e_1 = 2m$ and $e_2 = 2n$ for some $m, n \in \mathbb{Z}$. According to our assumption:
\[ 500 = o + e_1 + e_2 \]
Substituting our definitions:
\[ 500 = (2k + 1) + 2m + 2n \]
Factoring out a $2$ from the terms with variables:
\[ 500 = 2(k + m + n) + 1 \]
Since $k, m, n \in \mathbb{Z}$, their sum $(k + m + n)$ is an integer. Let $j = k + m + n$, where $j \in \mathbb{Z}$. Then $500 = 2j + 1$. By definition, any number of the form $2j + 1$ is an odd integer. This implies that $500$ is an odd integer. However, $500 = 2(250)$, which means $500$ is an even integer. We have reached a contradiction: $500$ cannot be both even and odd. Therefore, our assumption is false.
\end{proof}

\textbf{Solution 7:}
\begin{proof}
We proceed by contradiction. Assume that $a$ and $b$ are odd integers, and assume, to the contrary, that $4 \mid (a^2 + b^2)$. Since $a$ and $b$ are odd, there exist integers $k, m$ such that $a = 2k + 1$ and $b = 2m + 1$. Since $4 \mid (a^2 + b^2)$, there exists an integer $c$ such that $a^2 + b^2 = 4c$. We substitute the expressions for $a$ and $b$:
\begin{align*}
a^2 + b^2 &= (2k + 1)^2 + (2m + 1)^2 \\
&= (4k^2 + 4k + 1) + (4m^2 + 4m + 1) \\
&= 4k^2 + 4k + 4m^2 + 4m + 2 \\
&= 4(k^2 + k + m^2 + m) + 2
\end{align*}
Equating this to $4c$:
\[ 4c = 4(k^2 + k + m^2 + m) + 2 \]
Divide the entire equation by $2$:
\[ 2c = 2(k^2 + k + m^2 + m) + 1 \]
Let $j = k^2 + k + m^2 + m$. Since $k, m \in \mathbb{Z}$, $j \in \mathbb{Z}$. The equation becomes:
\[ 2c = 2j + 1 \]
The left side, $2c$, is an even integer. The right side, $2j + 1$, is an odd integer. An integer cannot be simultaneously even and odd. This is a contradiction. Thus, $4 \nmid (a^2 + b^2)$.
\end{proof}

\textbf{Solution 8:}
\begin{proof}
We proceed by contradiction. Assume, to the contrary, that there exist integers $x, y$ such that $x^2 - 4y = 2$. Rearranging the equation yields:
\[ x^2 = 4y + 2 = 2(2y + 1) \]
Since $y \in \mathbb{Z}$, $(2y + 1)$ is an integer. This implies $x^2$ is a multiple of $2$, so $x^2$ is even. If the square of an integer is even, the integer itself must be even. Let $x = 2k$ for some $k \in \mathbb{Z}$. Substitute this back into the equation:
\[ (2k)^2 = 2(2y + 1) \]
\[ 4k^2 = 4y + 2 \]
Divide the entire equation by $2$:
\[ 2k^2 = 2y + 1 \]
The left side, $2k^2 = 2(k^2)$, is an even integer. The right side, $2y + 1$, is an odd integer. This is a contradiction. Therefore, no such integer solutions exist.
\end{proof}

\textbf{Solution 9:}
\begin{proof}
We proceed by contradiction. Assume that the statement is false. By De Morgan's Laws, the negation of ``$a \nmid b$ or $a \nmid (b + a - 1)$'' is ``$a \mid b$ and $a \mid (b + a - 1)$''. Assume there exist integers $a \ge 2$ and $b$ such that $a \mid b$ and $a \mid (b + a - 1)$. By the definition of divisibility, there exist integers $k$ and $m$ such that:
\[ b = ak \]
\[ b + a - 1 = am \]
Substitute $b = ak$ into the second equation:
\[ ak + a - 1 = am \]
Rearranging the terms to isolate $1$:
\[ am - ak - a = -1 \implies a(k - m + 1) = 1 \]
Since $k, m \in \mathbb{Z}$, their combination $(k - m + 1)$ is an integer. Let $j = k - m + 1$, where $j \in \mathbb{Z}$. We have $a \cdot j = 1$. The only integer divisors of $1$ are $1$ and $-1$. Therefore, $a$ must be either $1$ or $-1$. However, our initial premise states that $a \ge 2$. This contradicts the deduction that $a \in \{1, -1\}$. Therefore, our assumption is false.
\end{proof}

\textbf{Solution 10:}
\begin{proof}
We proceed by contradiction. Assume, to the contrary, that there exists an integer $n$ such that $n^2$ is even, but $n$ is not even. If $n$ is not even, it must be odd. By definition of an odd integer, there exists an integer $k$ such that $n = 2k + 1$. We square $n$:
\[ n^2 = (2k + 1)^2 = 4k^2 + 4k + 1 = 2(2k^2 + 2k) + 1 \]
Since $k \in \mathbb{Z}$, let $m = 2k^2 + 2k$, where $m \in \mathbb{Z}$. Then $n^2 = 2m + 1$. This implies that $n^2$ is an odd integer. However, our initial premise stated that $n^2$ is even. A number cannot be both even and odd. This is a contradiction. Therefore, if $n^2$ is even, $n$ must be even.
\end{proof}

\textbf{Solution 11:}
\begin{proof}
We proceed by contradiction. Assume, to the contrary, that there exist three consecutive integers whose sum is $100$. Let the integers be $n, n+1, n+2$ for some $n \in \mathbb{Z}$. Their sum is:
\[ n + (n + 1) + (n + 2) = 100 \]
\[ 3n + 3 = 100 \]
\[ 3(n + 1) = 100 \]
Divide both sides by $3$:
\[ n + 1 = \frac{100}{3} \]
Since $n \in \mathbb{Z}$, $n+1$ must also be an integer. However, $\frac{100}{3}$ is not an integer because $100$ is not divisible by $3$ (remainder is $1$). This is a contradiction. Therefore, the sum of three consecutive integers cannot be $100$.
\end{proof}

\subsection*{Solutions to Part C: Rationality Bounds}

\textbf{Solution 12:}
\begin{proof}
We proceed by contradiction. Assume, to the contrary, that there exists a largest negative rational number. Let us call this number $r$. Since $r$ is a negative rational number, $r \in \mathbb{Q}$ and $r < 0$. Consider the number $\frac{r}{2}$. Since $r \in \mathbb{Q}$, there exist integers $p, q$ ($q \neq 0$) such that $r = \frac{p}{q}$. Then $\frac{r}{2} = \frac{p}{2q}$. Since $p, 2q \in \mathbb{Z}$ and $2q \neq 0$, it follows that $\frac{r}{2} \in \mathbb{Q}$. Since $r < 0$, dividing by the positive integer $2$ preserves the inequality, yielding $\frac{r}{2} < 0$. Thus, $\frac{r}{2}$ is a negative rational number. Furthermore, since $r < 0$, dividing by $2$ gives a number closer to zero. Formally, $r < \frac{r}{2} < 0$. This means $\frac{r}{2}$ is a negative rational number that is strictly greater than $r$. This contradicts our initial assumption that $r$ is the largest negative rational number.
\end{proof}

\textbf{Solution 13:}
\begin{proof}
We proceed by contradiction. Assume, to the contrary, that there exists a smallest positive irrational number, say $x$. Since $x$ is positive, $x > 0$. Consider the number $\frac{x}{2}$. Since $x > 0$, it follows that $\frac{x}{2} > 0$. We must show that $\frac{x}{2}$ is irrational. Assume for the sake of contradiction that $\frac{x}{2}$ is rational. Then $\frac{x}{2} = \frac{p}{q}$ for some $p, q \in \mathbb{Z}, q \neq 0$. This implies $x = \frac{2p}{q}$. Since $2p, q \in \mathbb{Z}$, this would mean $x$ is rational, which contradicts the fact that $x$ is irrational. Thus, $\frac{x}{2}$ must be irrational. Since $x > 0$, dividing by $2$ yields $\frac{x}{2} < x$. We have found a positive irrational number, $\frac{x}{2}$, which is strictly less than $x$. This contradicts the assumption that $x$ is the smallest positive irrational number.
\end{proof}

\textbf{Solution 14:}
\begin{proof}
We proceed by contradiction. Assume, to the contrary, that there exists a rational number $q$ and an irrational number $i$ such that their sum $S = q + i$ is a rational number. Since $S$ and $q$ are rational, they can be expressed as $S = \frac{a}{b}$ and $q = \frac{c}{d}$ for integers $a, b, c, d$ with $b \neq 0, d \neq 0$. We can isolate $i$:
\[ i = S - q = \frac{a}{b} - \frac{c}{d} = \frac{ad - bc}{bd} \]
Since $a, b, c, d$ are integers, $(ad - bc)$ is an integer and $bd$ is a non-zero integer. This implies that $i$ can be written as the ratio of two integers, meaning $i$ is a rational number. This contradicts our initial premise that $i$ is an irrational number. Therefore, the sum must be irrational.
\end{proof}

\subsection*{Solutions to Part D: De Morgan's Laws \& Sets}

\textbf{Solution 15:}
\begin{proof}
We proceed by contradiction. Assume that the statement is false. The negation of the implication $P \implies (Q \lor R)$ is $P \land \sim(Q \lor R)$, which by De Morgan's Laws is $P \land (\sim Q \land \sim R)$. Therefore, we assume that $x + y > 10$ AND $x \le 5$ AND $y \le 5$. 
From our assumptions, we can add the two inequalities $x \le 5$ and $y \le 5$:
\[ x + y \le 5 + 5 \implies x + y \le 10 \]
However, our initial premise states that $x + y > 10$. We have derived that $x + y \le 10$ and $x + y > 10$, which is a logical contradiction. Therefore, our assumption is false, and it must be true that $x > 5$ or $y > 5$.
\end{proof}

\textbf{Solution 16:}
\begin{proof}
We proceed by contradiction. Assume that the premise $A \cap B = \emptyset$ is true, but the conclusion $A \subseteq B^c$ is false. By the definition of subsets, if $A \not\subseteq B^c$, there must exist at least one element $x \in A$ such that $x \notin B^c$. 
By the definition of a complement, if $x \notin B^c$, then $x \in B$. 
Since we established $x \in A$ and $x \in B$, it follows by the definition of intersection that $x \in A \cap B$. 
This implies that $A \cap B$ contains at least one element (namely, $x$), meaning $A \cap B \neq \emptyset$. This directly contradicts our initial premise that $A \cap B = \emptyset$. Therefore, $A \subseteq B^c$.
\end{proof}

\subsection*{Solutions to Part E: Infinite Descent}

\textbf{Solution 17:}
\begin{proof}
Assume, to the contrary, that $\sqrt{3}$ is rational. Then there exist positive integers $p, q$ such that $\sqrt{3} = \frac{p}{q}$. Squaring both sides gives $3 = \frac{p^2}{q^2}$, so $p^2 = 3q^2$. Since $p^2$ is a multiple of $3$, $p$ must be a multiple of $3$ (since $3$ is prime). Let $p = 3p_1$ for some positive integer $p_1 < p$. Substitute $p$ into the equation:
\[ (3p_1)^2 = 3q^2 \implies 9p_1^2 = 3q^2 \implies q^2 = 3p_1^2 \]
This implies $q^2$ is a multiple of $3$, so $q$ must be a multiple of $3$. Let $q = 3q_1$ for some positive integer $q_1 < q$. We now have a new pair of positive integers $(p_1, q_1)$ such that $p_1^2 = 3q_1^2$, meaning $\sqrt{3} = \frac{p_1}{q_1}$. By repeating this process, we generate an infinite sequence of strictly decreasing positive integers:
\[ p > p_1 > p_2 > p_3 > \dots > 0 \]
This contradicts the Well-Ordering Principle of the natural numbers (an infinite strictly decreasing sequence cannot exist in $\mathbb{N}$). Therefore, $\sqrt{3}$ is irrational.
\end{proof}

\textbf{Solution 18:}
\begin{proof}
Assume, to the contrary, that there exists a non-zero integer solution $(x, y, z)$. We can assume without loss of generality that $x, y, z$ are positive. We have:
\[ x^3 + 3y^3 = 9z^3 \]
Rearranging gives $x^3 = 9z^3 - 3y^3 = 3(3z^3 - y^3)$. This implies $3 \mid x^3$. Since $3$ is prime, $3 \mid x$. Let $x = 3x_1$ for some positive integer $x_1 < x$. Substitute this back:
\[ (3x_1)^3 + 3y^3 = 9z^3 \implies 27x_1^3 + 3y^3 = 9z^3 \]
Divide by $3$:
\[ 9x_1^3 + y^3 = 3z^3 \implies y^3 = 3(z^3 - 3x_1^3) \]
This implies $3 \mid y^3$, so $3 \mid y$. Let $y = 3y_1$ for some positive integer $y_1 < y$. Substitute this back:
\[ 9x_1^3 + (3y_1)^3 = 3z^3 \implies 9x_1^3 + 27y_1^3 = 3z^3 \]
Divide by $3$:
\[ 3x_1^3 + 9y_1^3 = z^3 \implies z^3 = 3(x_1^3 + 3y_1^3) \]
This implies $3 \mid z^3$, so $3 \mid z$. Let $z = 3z_1$ for some positive integer $z_1 < z$. We have found a new, strictly smaller positive integer solution $(x_1, y_1, z_1)$. This process can be repeated infinitely, creating an infinite descent of positive integers, which violates the Well-Ordering Principle. Thus, no non-zero integer solutions exist.
\end{proof}

\textbf{Solution 19:}
\begin{proof}
Assume, to the contrary, that there exists a positive integer solution $(x, y, z)$ to $x^2 + y^2 = 3z^2$. We analyze the equation modulo $3$. The squares of integers modulo $3$ can only be $0$ or $1$ (since $0^2 \equiv 0$, $1^2 \equiv 1$, $2^2 \equiv 1 \pmod{3}$). 
The right side is $3z^2 \equiv 0 \pmod{3}$. 
For the left side $x^2 + y^2$ to be congruent to $0 \pmod{3}$, both $x^2$ and $y^2$ must be congruent to $0 \pmod{3}$ (since $1+0=1$ and $1+1=2$). 
Thus, $3 \mid x^2$ and $3 \mid y^2$, which implies $3 \mid x$ and $3 \mid y$. Let $x = 3x_1$ and $y = 3y_1$ for positive integers $x_1 < x$ and $y_1 < y$. Substitute back:
\[ (3x_1)^2 + (3y_1)^2 = 3z^2 \implies 9x_1^2 + 9y_1^2 = 3z^2 \implies 3(x_1^2 + y_1^2) = z^2 \]
This implies $3 \mid z^2$, so $3 \mid z$. Let $z = 3z_1$ for $z_1 < z$. We have a new, strictly smaller positive integer solution $(x_1, y_1, z_1)$. This creates an infinite descent, violating the Well-Ordering Principle. Thus, no positive integer solutions exist.
\end{proof}

\subsection*{Solutions to Part F: Probability \& Abstract Modeling}

\textbf{Solution 20:}
\begin{proof}
Assume, to the contrary, that $\mathbb{P}(B) \neq 1$. Since probabilities are bounded by $1$, this means $\mathbb{P}(B) < 1$. By the axioms of probability, for any two events $A$ and $B$:
\[ \mathbb{P}(A \cup B) = \mathbb{P}(A) + \mathbb{P}(B) - \mathbb{P}(A \cap B) \]
We are given $\mathbb{P}(A) = 0$. Since $A \cap B \subseteq A$, monotonicity implies $0 \le \mathbb{P}(A \cap B) \le \mathbb{P}(A) = 0$, so $\mathbb{P}(A \cap B) = 0$. Substituting these values into the formula yields:
\[ \mathbb{P}(A \cup B) = 0 + \mathbb{P}(B) - 0 = \mathbb{P}(B) \]
Since we assumed $\mathbb{P}(B) < 1$, it follows that $\mathbb{P}(A \cup B) < 1$. However, our initial premise states that $\mathbb{P}(A \cup B) = 1$. This is a contradiction ($1 < 1$). Therefore, $\mathbb{P}(B) = 1$.
\end{proof}

\textbf{Solution 21:}
\begin{proof}
Assume, to the contrary, that $A \cap B = \emptyset$. This means $A$ and $B$ are mutually exclusive events. By the third axiom of probability (countable additivity for disjoint sets):
\[ \mathbb{P}(A \cup B) = \mathbb{P}(A) + \mathbb{P}(B) \]
By the first axiom of probability, the probability of any event is bounded above by $1$, so $\mathbb{P}(A \cup B) \le 1$. Substituting our sum yields:
\[ \mathbb{P}(A) + \mathbb{P}(B) \le 1 \]
However, our initial premise states that $\mathbb{P}(A) + \mathbb{P}(B) > 1$. We have derived a contradiction. Therefore, our assumption that $A \cap B = \emptyset$ is false, meaning $A \cap B \neq \emptyset$.
\end{proof}

\textbf{Solution 22:}
\begin{proof}
Assume, to the contrary, that an event $A$ is independent of itself, but $\mathbb{P}(A) \neq 0$ and $\mathbb{P}(A) \neq 1$. By the definition of independence, two events $X$ and $Y$ are independent if $\mathbb{P}(X \cap Y) = \mathbb{P}(X)\mathbb{P}(Y)$. Since $A$ is independent of itself, we set $X = A$ and $Y = A$:
\[ \mathbb{P}(A \cap A) = \mathbb{P}(A)\mathbb{P}(A) \]
Since $A \cap A = A$, this simplifies to:
\[ \mathbb{P}(A) = (\mathbb{P}(A))^2 \]
Let $x = \mathbb{P}(A)$. The equation is $x = x^2$, which rearranges to $x^2 - x = 0$, or $x(x - 1) = 0$. The only real solutions to this equation are $x = 0$ and $x = 1$. This means $\mathbb{P}(A)$ must be $0$ or $1$. This contradicts our assumption that $\mathbb{P}(A) \neq 0$ and $\mathbb{P}(A) \neq 1$. Therefore, $\mathbb{P}(A)$ must be $0$ or $1$.
\end{proof}

\textbf{Solution 23:}
\begin{proof}
Assume, to the contrary, that the market contains no arbitrage. By the Fundamental Theorem of Asset Pricing, there must exist a risk-neutral probability measure $\mathbb{Q}$ such that the discounted asset price is a martingale: $S_0 = \mathbb{E}^\mathbb{Q}\left[\frac{S_1}{1+r}\right]$. Let $q = \mathbb{Q}(S_1 = 110)$. Then $\mathbb{Q}(S_1 = 120) = 1 - q$. Since $\mathbb{Q}$ is a probability measure, $0 \le q \le 1$. We set up the pricing equation:
\[ 100 = \frac{110q + 120(1 - q)}{1.05} \]
Multiply by $1.05$:
\[ 105 = 110q + 120 - 120q \]
\[ 105 = 120 - 10q \]
\[ 10q = 15 \implies q = 1.5 \]
We have derived that the probability $q = 1.5$. However, by the axioms of probability, $q \le 1$. This is a contradiction ($1.5 \le 1$). Therefore, no such measure $\mathbb{Q}$ exists, meaning the market must contain an arbitrage opportunity.
\end{proof}

\textbf{Solution 24:}
\begin{proof}
Assume, to the contrary, that the market contains no arbitrage. We construct a trading strategy: at time $t=0$, buy portfolio $A$ and short-sell (sell) portfolio $B$. The initial cash flow is $V_0 = V_0(B) - V_0(A)$. Since $V_0(A) < V_0(B)$, $V_0 > 0$, meaning we receive strictly positive cash today. We invest this cash in the risk-free bond.
At time $T$, the value of our holdings is the payoff of $A$ minus the payoff of $B$, plus the bond value:
\[ V_T = V_T(A) - V_T(B) + V_0(1+r)^T \]
Since $V_T(A) = V_T(B)$ almost surely, $V_T(A) - V_T(B) = 0$. Thus, $V_T = V_0(1+r)^T$. Since $V_0 > 0$ and $r \ge 0$, $V_T > 0$ almost surely. 
We have constructed a portfolio with zero initial cost (actually, positive initial cash extracted), no risk of loss ($\mathbb{P}(V_T \ge 0) = 1$), and a strict profit ($\mathbb{P}(V_T > 0) = 1 > 0$). This perfectly satisfies the definition of an arbitrage. This contradicts our assumption that the market contains no arbitrage.
\end{proof}

\textbf{Solution 25:}
\begin{proof}
Assume, to the contrary, that an arbitrage strategy $(x, y)$ exists AND an equivalent martingale measure $\mathbb{Q}$ exists. By definition of arbitrage, the initial portfolio value is $V_0 = 0$, and at time $t=1$, $\mathbb{P}(V_1 \ge 0) = 1$ and $\mathbb{P}(V_1 > 0) > 0$. 
Since $\mathbb{Q}$ is equivalent to $\mathbb{P}$, they agree on probability zero events. Thus, $\mathbb{Q}(V_1 \ge 0) = 1$ and $\mathbb{Q}(V_1 > 0) > 0$. 
The expectation of a non-negative random variable that is strictly positive with non-zero probability must be strictly positive. Therefore, $\mathbb{E}^\mathbb{Q}[V_1] > 0$.
However, since $\mathbb{Q}$ is a martingale measure, the discounted value of any self-financing portfolio is a martingale. Thus, $V_0 = \mathbb{E}^\mathbb{Q}\left[\frac{V_1}{1+r}\right]$. Since $V_0 = 0$, we have:
\[ 0 = \frac{1}{1+r} \mathbb{E}^\mathbb{Q}[V_1] \implies \mathbb{E}^\mathbb{Q}[V_1] = 0 \]
We have derived that $\mathbb{E}^\mathbb{Q}[V_1] > 0$ and $\mathbb{E}^\mathbb{Q}[V_1] = 0$. This is a contradiction. Therefore, $\mathbb{Q}$ cannot exist.
\end{proof}

\textbf{Solution 26:}
\begin{proof}
Assume, to the contrary, that a hypothesis class $\mathcal{H}$ with infinite VC-dimension ($\text{VC}(\mathcal{H}) = \infty$) can be PAC learned by some algorithm $\mathcal{A}$. By the definition of PAC learning, for any $\epsilon \in (0, 1)$ and $\delta \in (0, 1)$, there exists a finite sample complexity $m_{\mathcal{H}}(\epsilon, \delta) < \infty$ such that $\mathcal{A}$ successfully learns $\mathcal{H}$.
By the Fundamental Theorem of Statistical Learning, the sample complexity is bounded below:
\[ m_{\mathcal{H}}(\epsilon, \delta) \ge C \frac{\text{VC}(\mathcal{H}) + \ln(1/\delta)}{\epsilon} \]
Substitute $\text{VC}(\mathcal{H}) = \infty$ into the inequality:
\[ m_{\mathcal{H}}(\epsilon, \delta) \ge \infty \]
This implies that the required sample size $m$ must be infinite. However, our assumption of PAC learnability requires $m$ to be finite. We have derived that $m$ is both finite and infinite, a contradiction. Therefore, $\mathcal{H}$ cannot be PAC learned.
\end{proof}

\subsection*{Solutions to Part G: Advanced Diagnostic Testing}

\textbf{Solution 27:}
\begin{proof}
\textbf{The Flaw:} The proof divides both sides of the equation by the term $(a - b)$. Since the proof explicitly defines $a = 1$ and $b = 1$, it follows that $a - b = 1 - 1 = 0$. Division by zero is undefined and strictly violates the field axioms of real numbers. Any logical deduction following a division by zero is invalid (Principle of Explosion). The contradiction is fabricated by an illegal algebraic operation.
\end{proof}

\textbf{Solution 28:}
\begin{proof}
\textbf{The Flaw:} The logical deduction ``if $k$ is even, $q$ is even'' is a conditional statement that is never proven to be triggered. The proof establishes $k^2 = q^2$, but it never proves that $k$ must be even. In reality, for $\sqrt{4} = \frac{2}{1}$, $p = 2, q = 1$, so $k = 1$. Since $k$ is odd, $q$ is odd, and they share no common factors. The contradiction is fabricated by an unproven ``if'' statement.
\end{proof}

\textbf{Solution 29:}
\begin{proof}
\textbf{The Flaw:} The Intermediate Value Theorem (IVT) strictly requires the function $f(x)$ to be continuous on the closed interval $[a, b]$. The function $f(x) = \frac{1}{x}$ has an infinite discontinuity at $x = 0$, which lies within the interval $[-1, 1]$. Therefore, the IVT cannot be applied to this function on this interval, and the derived contradiction is invalid.
\end{proof}

\textbf{Solution 30:}
\begin{proof}
\textbf{The Flaw:} The proof assumes that if $(x, y)$ are positive integers, then $(x - 1, y - 1)$ are also positive integers. This is false if $x = 1$ or $y = 1$. If $x = 1$, then $x - 1 = 0$, which is not a positive integer, meaning $(x - 1, y - 1) \notin S$. The infinite descent fails at the base case, rendering the contradiction invalid.
\end{proof}

\newpage

% ==========================================
% SECTION 4: CONCEPT CHECKLIST (REPEATED)
% ==========================================
\section{Concept Checklist \& Question Mapping (Repeated)}
\label{sec:checklist_end}

\begin{itemize}
    \item \textbf{C1: Universal Falsification (Finite Domains):} Exhaustive testing to find edge cases where the hypothesis holds but the conclusion fails. \textit{(Mapped to Problems 1, 2)}
    \item \textbf{C2: Universal Falsification (Infinite Domains):} Analyzing algebraic structures, inequalities, and domain restrictions (e.g., logarithms, roots) to construct counterexamples. \textit{(Mapped to Problems 3, 4, 5)}
    \item \textbf{C3: Impossibility Proofs (Parity Forcing):} Mapping complex polynomial expressions to the binary space of $\{0, 1\} \pmod{2}$ to expose contradictions. \textit{(Mapped to Problems 6, 7, 8)}
    \item \textbf{C4: Impossibility Proofs (Divisibility \& Primes):} Forcing violations of the fundamental theorem of arithmetic or basic divisibility axioms. \textit{(Mapped to Problems 9, 10, 11)}
    \item \textbf{C5: Rationality Bounds \& Density:} Exploiting the density of $\mathbb{Q}$ and $\mathbb{R} \setminus \mathbb{Q}$ using the ``Divide by 2'' trick to shatter extremum assumptions. \textit{(Mapped to Problems 12, 13, 14)}
    \item \textbf{C6: De Morgan's Laws on Implications:} Negating disjunctive conclusions ($\sim(P \implies Q \lor R) \equiv P \land \sim Q \land \sim R$) to form concrete simultaneous equations. \textit{(Mapped to Problems 15, 16)}
    \item \textbf{C7: Fermat's Method of Infinite Descent:} Constructing an infinite strictly decreasing sequence of positive integers to violate the Well-Ordering Principle of $\mathbb{N}$. \textit{(Mapped to Problems 17, 18, 19)}
    \item \textbf{C8: Set-Theoretic \& Probability Foundations:} Utilizing mutually exclusive/exhaustive sets and probability measure axioms ($\mathbb{P}(A)=1$) to prove logical necessities. \textit{(Mapped to Problems 20, 21, 22)}
    \item \textbf{C9: Abstract Real-World Modeling (Quant Finance):} Constructing self-financing portfolios to contradict the existence of an Equivalent Martingale Measure ($\mathbb{Q}$) in Arbitrage Pricing Theory. \textit{(Mapped to Problems 23, 24, 25)}
    \item \textbf{C10: Abstract Real-World Modeling (Machine Learning):} Proving PAC learning lower bounds by assuming learnability for infinite VC-dimension classes and deriving contradictions. \textit{(Mapped to Problems 26)}
    \item \textbf{C11: Advanced Diagnostic Testing (Find the Flaw):} Identifying subtle logical fallacies, such as unproven conditionals, division by zero, or invalid theorem applications (e.g., IVT on discontinuous functions). \textit{(Mapped to Problems 27, 28, 29, 30)}
\end{itemize}

\end{document}